                                  IMC 2020 Online
                                  Day 1, July 26, 2020




Problem 1. Let      n be a positive integer. Compute the number of words w (nite sequences of
letters) that satisfy all the following three properties:

 (1) w consists of n letters, all of them are from the alphabet {a, b, c, d};

 (2) w contains an even number of letters a;

 (3) w contains an even number of letters b.

(For example, for n = 2 there are 6 such words: aa, bb, cc, dd, cd and dc.)
                                                Armend Sh. Shabani, University of Prishtina

Solution 1.    Let N = {1, 2, . . . , n}. Consider a word w that satises the conditions and let
A, B, C, D ⊂ N be the sets of positions of letters a, b, c and d in w, respectively. By the
denition of the words we have A t B t C t D = N . The sets A and B are constrained to have
even sizes.
    In order to construct all suitable words w, choose the set S = A ∪ B rst; by the conditions,
|S| = |A| + |B| must be even. It is well-known that an n-element set (with n ≥ 1) has 2n−1
even subsets, so there are 2n−1 possibilities for S .
    If S = ∅ then we can choose C ⊂ N arbitrarily, and then the set D = S \ C is determined
D uniquely. Since N has 2n subsets, we have 2n options for set C and therefore 2n suitable
words w with S = ∅.
    Otherwise, if k = |S| > 0, we have to choose an arbitrary subset C of N \ S and an even
subset A of S ; then D = (N \ S) \ C and B = S \ A are determined and |B| = |S| − |A| will
automatically be even. We have 2n−k choices for C and 2k−1 independent choices for A; so for
each nonempty even S we have 2n−k · 2k−1 = 2n−1 suitable words.
    The number of nonempty even sets S is 2n−1 − 1, so in total, the number of words satisfying
the conditions is
                             1 · 2n + (2n−1 − 1) · 2n−1 = 4n−1 + 2n−1 .
   Solution 2.   Let an denote the number of words of length n over A = {a, b, c, d} such
that a and b appear even number of times. Further, we dene the following sequences for the
number of words of length n, all over A.

   • bn - the number of words with an odd number of a's and even number of b's

   • cn - the number of words with even number of a's and an odd number of b's

   • dn - the number of words with an odd number of a's and an odd number of b's

We will call them A-words, B-words, C-words and D-words, respectively.
It is clear that a1 = 2 and that

                                    an + bn + cn + dn = 4n .

First, we nd a recurrence relation for an . If an A-word of length n begins with c or d, it can
be followed by any A-word of length n − 1, contributing with 2an−1 . If an A-word of length
n begins with a, it can be followed by any word of length n − 1 that contains an odd number

                                                1
of a's and even number of b's, thus contributing with bn−1 . If an A-word of length n begins
with b, it can be followed by any word of length n − 1 that contains even number of a's and
an odd number of b's, thus contributing with cn−1 . Therefore we have the following recurrence
relation:
                                  an = 2an−1 + bn−1 + cn−1 .                               (1)
Next, we nd a recurrence relation for bn .
If a B-word of length n begins with c or d, it can be followed by any B-word of length n − 1,
contributing with 2bn−1 . If a B-word of length n begins with a, it can be followed by any word of
length n−1 that contains even number of a's and even number of b's, contributing with an−1 . If a
B-word of length n begins with b, it can be followed by any word of length n−1 that contains an
odd number of a's and an odd number of b's, contributing with dn−1 = 4n−1 −an−1 −bn−1 −cn−1 .
Therefore we have the following recurrence relation:

                                    bn = bn−1 + 4n−1 − cn−1 .                                 (2)

   Now observe that bk = ck for all k , since simultaneously replacing a's to b's and vice versa
we get a C -word from a B -word. Therefore (2) yields bn = 4n−1 . Now (1) yields

                                     an = 2 · an−1 + 2 · 4n−2 .

Solving the last recurrence relation (for example, diving by 2n we get xn := an 2−n satises
xn − xn−1 = 2n−3 , and it remains to sum up consecutive powers of 2) we get

                                        an = 2n−1 + 4n−1 .
   Solution 3. Consider the sum

         (a + b + c + d)n + (−a − b + c + d)n + (−a + b + c + d)n + (a − b + c + d)n
                                                                                     .        (∗)
                                              4
Expanding the parentheses as

              (a + b + c + d)n = (a + b + c + d)(a + b + c + d) . . . (a + b + c + d),

we get a sum of products x1 . . . xn , xi ∈ {a, b, c, d}, naturally corresponding to the words of
length n over the alphabet {a, b, c, d}. Consider the other terms in the numerator similarly.
   If a word x1 . . . xn contains A, B, C, D letters a,b,c and d respectively, we get aA bB cC dD
with the coecient
                                                                      (
    1 + (−1)A+B + (−1)A + (−1)B         (1 + (−1)A )(1 + (−1)B )        1, if A and B are even
                                      =                            =
                    4                                 4                 0, otherwise.

Hence, by substituting a = b = c = d = 1 in (∗) we get the answer (4n + 2n+1 )/4 = 4n−1 + 2n−1 .




                                                 2
Problem 2. Let    A and B be n × n real matrices such that

                                      rk(AB − BA + I) = 1

where I is the n × n identity matrix.
    Prove that
                                                          1
                          trace(ABAB) − trace(A2 B2 ) = n(n − 1).
                                                          2
(rk(M ) denotes the rank of matrix M , i.e., the maximum number of linearly independent
columns in M . trace(M) denotes the trace of M , that is the sum of diagonal elements in M .)
                              Rustam Turdibaev, V. I. Romanovskiy Institute of Mathematics

Solution. Let   X = AB − BA. The rst important observation is that

   trace(X2 ) = trace(ABAB − ABBA − BAAB + BABA) = 2trace(ABAB) − 2trace(A2 B2 )

using that the trace is cyclic. So we need to prove that trace(X2 ) = n(n − 1).
   By assumption, X + I has rank one, so we can write X + I = v t w for two vectors v, w. So

                  X 2 = (v t w − I)2 = I − 2v t w + v t wv t w = I + (wv t − 2)v t w.

Now by denition of X we have trace(X) = 0 and hence wv t = trace(wvt ) = trace(vt w) = n so
that indeed
                          trace(X2 ) = n + (n − 2)n = n(n − 1).
   An alternative way to use the rank one condition is via eigenvalues: Since X + I has rank
one, it has eigenvalue 0 with multiplicity n − 1. So X has eigenvalue −1 with multiplicity n − 1.
Since trace(X) = 0 the remaining eigenvalue of X must be n − 1. Hence

                         trace(X2 ) = (n − 1)2 + (n − 1) · 12 = n(n − 1).




                                                  3
Problem 3.     Let d ≥ 2 be an integer. Prove that there exists a constant C(d) such that the
following holds: For any convex polytope K ⊂ Rd , which is symmetric about the origin, and
any ε ∈ (0, 1), there exists a convex polytope L ⊂ Rd with at most C(d)ε1−d vertices such that

                                        (1 − ε)K ⊆ L ⊆ K.

(For a real α, a set T ⊂ Rd with nonempty interior is a convex polytope
                                                                   P with at most αP vertices, if
T is a convex hull of a set X ⊂ R of at most α points, i.e., T = { x∈X tx x | tx ≥ 0, x∈X tx =
                                  d

1}. For a real λ, put λK = {λx | x ∈ K}. A set T ⊂ Rd is symmetric about the origin if
(−1)T = T .)
                                                 Fedor Petrov, St. Petersburg State University

Solution [in elementary terms] Let      {p1 , . . . , pm } be an inclusion-maximal collection of points
on the boundary ∂K of K such that the homothetic copies Ki := pi + 2ε K have disjoint interiors.
We claim that the convex hull L := conv{p1 , . . . , pm } satises all the conditions.
   First, note that by convexity of K we have aK + bK = (a + b)K for a, b > 0. It follows that
Ki ⊂ (1 + 2ε )K . On the other hand, if k ∈ K , a > 0 and and ak ∈ Ki , then
                                       ε        ε        ε
                              pi ∈ ak − K = ak + K ⊂ (a + )K,
                                       2        2        2
and since pi is a boundary point of K , we get a + 2ε ⩾ 1, a ⩾ 1 − 2ε . It means that all Ki lie
between (1 − 2ε )K and (1 + 2ε )K . Since their interiors are disjoint, by the volume counting we
obtain                           ε d          ε d       ε d
                             m          ⩽ 1+         − 1−         ⩽ (3/2)d ε
                                  2              2           2
(since F (ε) = (1 + 2ε )d − (1 − 2ε )d is a polynomial in ε without constant term with non-negative
coecients which sum up to (3/2)d − (1/2)d ), therefore m ⩽ 3d ε1−d .
    It is clear that L ⊆ K , so it remains to prove that (1 − ε)K ⊆ L. Assume the contrary:
there exists a point p ∈ (1 − ε)K \ L. Separate p from L by a hyperplane: Choose a linear
functional ` such that `(p) > maxx∈L `(x) = maxi `(pi ). Choose x ∈ K such that `(x) =: a is
maximal possible. Note that by our construction x + 2ε K has a common point with some Ki :
there exists a point z ∈ (x + 2ε K) ∩ (pi + 2ε K). We have
                                          ε                 ε
                                  `(pi ) + a ⩾ `(z) ⩾ `(x) − a,
                                          2                 2
and therefore `(pi ) ⩾ a(1 − ε). Since p ∈ (1 − ε)K , we obtain `(p) ⩽ a(1 − ε). A contradiction.
    Solution [in the language of Banach spaces] Equip R with the norm k · k, whose unit
                                                                  d

ball is K , call this Banach space V . Choose an inclusion maximal set X ⊂ ∂K whose pairwise
distances are ≥ ε. Put L = convX .
    The inclusion L ⊆ K follows from the convexity of K . If the inclusion (1−ε)K ⊆ L fails then
the HahnBanach theorem provides a unit linear functional λ ∈ V ∗ such that max{λ(L)} =
max{λX} ≤ 1 − ε. Then the point x ∈ K , where the maximum max{λ(K)} = 1 is attained
(thanks to the nite dimension and compactness) is in ∂K and, as λ witnesses, at distance ≥ ε
from all other points of L and X , contradicting the inclusion-maximality of X .
    The upper bound for the cardinality |X| is obtained by noting that the ε/2 balls centered
at the points of X are pairwise disjoint and lie in the dierence of balls (1 + ε/2)K \ (1 − ε/2)K ,
whose volume is (1 + ε/2)d − (1 − ε/2)d volK , the volume of each of the small balls being
εd /2d volK . Hence
                                      (2 + ε)d − (2 − ε)d
                               |X| ≤                       = O(ε1−d ).
                                               εd


                                                  4
Problem 4. A polynomial     p with real coecients satises the equation p(x + 1) − p(x) = x100
for all x ∈ R. Prove that p(1 − t) ⩾ p(t) for 0 ⩽ t ⩽ 1/2.
                                                 Daniil Klyuev, St. Petersburg State University

Solution 1.     Denote h(z) = p(1 − z̄) − p(z) for complex z . For t ∈ R we have h(it) =
p(1 + it) − p(it) = t100 , h(1/2 + it) = 0.
   If p(z) = cn z n + . . . + c0 , cn 6= 0, we have

           h(a + it) = p((1 − a) + it) − p(a + it) = (1 − 2a) ncn in−1 tn−1 + Q(t, a)
                                                                                     


for some polynomial Q having degree at most n − 2 with respect to the variable t. Substituting
a = 0 we get n = 101, cn = 1/101.
    Next, for large |t| we see that <(h(a + it)) > 0 for 0 ⩽ a < 1/2.
    Therefore by Maximum Principle for the harmonic function <h and the rectangle [0, 1/2] ×
[−N, N ] for large enough N we conclude that <h is non-negative in this rectangle, in particular
on [0, 1/2], as we need.      Pm
    Solution 2. Let p(x) =       j=0 aj x . Then
                                         j


                  m                                                                             
                  X
                                   j      j
                                                                            m−1    m m−2
p(x+1)−p(x) =            aj (x + 1) − x           = a1 +a2 (2x+1)+· · ·+am mx    +     x  + ··· + 1 .
                   j=0
                                                                                    2

This implies that m = 101, mam = 1 so a101 = 101       1
                                                           , (m − 1)am−1 + am m            = 0 so a100 = − 12
                                                                                        
                                                                                     2
etc. For j ⩾ 1 aj is uniquely dened, a0 may be chosen arbitrarily.
     The equality p2n ( 21 ) = 0 holds because 0 = p2n ( 21 ) + p2n (1 − 12 ) = 2p2n ( 12 ). Let n ⩾ 1 be an
integer and let pn be a polynomial such that pn (x + 1) − pn (x) = xn for all x and pn (0) = 0 =
pn (1). The above considerations prove the uniqueness of pn . We have p1 (x) = 21 x2 − 12 x. Also
p0n (x + 1) − p0n (x) = nxn−1 = n (pn−1 (x + 1) − pn−1 (x)). Therefore p0n (x) = npn−1 (x) + cn−1 for
a properly chosen constant cn−1 . We shall prove that

(1) p2n−1 (x) − p2n−1 (1 − x) = 0, p2n (x) + p2n (1 − x) = 0, c2n = 0, p002n (x) = 2n(2n − 1)p2n−2 (x)

for n = 1, 2, . . . and for all x. Simple computation shows that p1 (x) − p1 (1 − x) = 0. We have
(p2 (x) + p2 (1 − x))0 = 2p1 (x) + c1 − (2p1 (1 − x) + c1 ) = 0 so the map x 7→ p2 (x) + p2 (1 − x) is
constant thus p2 (x) + p2 (1 − x) = p2 (0) + p2 (1 − 0) = 0. If the rst two equalities hold for some
n then (p2n+1 (x) − p2n+1 (1 − x))0 = (2n + 1)p2n (x) + c2n + (p2n (1 − x) + c2n ) = 2c2n so there
exists b ∈ R such that p2n+1 (x) − p2n+1 (1 − x) = 2c2n x + b for all x. p2n+1 (0) − p2n+1 (1 − 0) = 0
and p2n+1 (1)−p2n+1 (1−1) = 0 so 2c2n = 0 = b. This proves that p2n+1 (x)−p2n+1 (1−x) = 0 for
all x. In a similar way we shall prove the second equality: (p2n+2 (x) + p2n+2 (1 − x))0 = (2n +
2)p02n+1 (x) + c2n+1 − (2n + 2) (p2n+1 (1 − x) + c2n+1 ) = 0 so the map x 7→ p2n+2 (x) + p2n+2 (1 − x)
is constant hence p2n+2 (x) + p2n+2 (1 − x) = p2n+2 (0) + p2n+2 (1 − 0) = 0 for all x. Now
p002n+2 (x) = ((2n + 2)p2n+1 (x) + c2n+1 )0 = (2n + 2)p02n+1 (x) = (2n + 2)((2n + 1)p2n (x) + c2n ) =
(2n + 2)(2n + 1)p2n (x). Since p02 (x) = 2p1 (x) + c1 = x2 − x + c1 we obtain p002 (x) = 2x − 1 < 0 for
x < 21 .The function p2 is strictly concave on [0, 21 ] and p2 (0) = 0 = p2 ( 21 ). Therefore p2 (x) > 0
for x ∈ (0, 21 ). This together with the equality p4 (x) = 12p2 (x) implies that p4 is strictly convex
on [0, 21 ] so in view of p4 (0) = 0 = p4 ( 12 ) we conclude that p4 (x) < 0 for x ∈ (0, 12 ). Easy
induction shows that for x ∈ (0, 21 ) one has p2n (x) > 0 for an odd n and p2n (x) < 0 for an even
n. If t ∈ (0, 21 ) then by (1) we get p100 (1 − t) − p100 (t) = −2p100 (t) > 0 as required.




                                                         5
